\begin{figure}[tb]
\centering
\begin{tikzpicture}[>=Latex, font=\small,
  box/.style={draw,rounded corners=2pt,minimum height=10mm,align=center},
  state/.style={align=center,font=\footnotesize}]
  \node[box,minimum width=19mm] (e) at (0,0) {$\mathsf U_E$};
  \node[box,minimum width=39mm] (f) at (4.8,0)
       {$\mathsf C$\,: one recursive call};
  \node[box,minimum width=19mm] (r) at (9.6,0) {$\mathsf U_R^*$};
  \draw[->] (-1.8,0) -- (e.west);
  \draw[->] (e.east) -- (f.west);
  \draw[->] (f.east) -- (r.west);
  \draw[->] (r.east) -- (11.4,0);
  \node[state] at (-1.45,.58) {$\iota x$};
  \node[state] at (2.05,.58) {$\jmath\mathcal E_Lx$};
  \node[state] at (7.55,.58) {$\jmath\mathcal R_LVx$};
  \node[state] at (11.02,.58) {$\iota Vx$};
  \node[state] at (0,-.92) {local encoding};
  \node[state] at (4.8,-.92) {order $k<d$; controlled by $z$\\
                           fresh workspace starts at zero};
  \node[state] at (9.6,-.92) {local decoding};
  \draw[decorate,decoration={brace,mirror,amplitude=4pt}]
       (1.22,-1.7) -- (8.38,-1.7);
  \node[state] at (4.8,-2.12)
       {comparison states lie in the fixed good subspace};
\end{tikzpicture}
\caption{One recursion level.  Labels above the arrows are the comparison
states used in Lemma~\ref{rec:step}; the intermediate states are approximate,
as quantified there.  The phase register supplies the control $z$ to the
single smaller synthesis circuit.  The local encoders contain their own
amplification rounds, while the smaller circuit occurs only once.  The map
$\iota$ initializes all work registers to zero, so the final comparison
also records their cleanup.}
\label{fig:recursion}
\end{figure}
